\section{Legal Foundations and Rationale}

Courts require civilians to invoke their constitutional protections clearly and precisely. Vague, hesitant, or incomplete speech can be interpreted as consent or waiver, leading to coerced statements or lost defenses. The Invocation of Rights addresses this vulnerability through a standardized, four‑line script, each line grounded in constitutional doctrine and Supreme Court precedent. Although the legal effect of each line varies by context, their collective function is to assert core rights across the three common encounter types—consensual, detained, and arrested—throughout all jurisdictions.

\subsection*{``I invoke my right to remain silent.''}
This line activates the Fifth Amendment’s protection against self‑incrimination. In \textit{Berghuis v.\ Thompkins} (2010) the Supreme Court found that a suspect who stayed silent for nearly three hours post‑Mirandization, then answered one question, had waived the right—because silence alone did not clearly invoke it. In \textit{Salinas v.\ Texas} (2013) a plurality held that pre‑custodial silence could be treated as evidence of guilt unless the right was explicitly invoked. Clear phrasing eliminates ambiguity, ensuring that courts recognize the speaker’s intent before or after custody attaches.

\subsection*{``I invoke my right to a lawyer.''}
This line invokes the right to counsel during custodial interrogation (Fifth Amendment) and after formal charges (Sixth). In \textit{Davis v.\ United States} (1994) the Court ruled that ``Maybe I should talk to a lawyer'' was too equivocal to require police to stop questioning. \textit{Edwards v.\ Arizona} (1981) held that once a clear request for counsel is made, police must cease interrogation unless the suspect re‑initiates. While the line does not compel police to stop questioning before custody attaches, it creates a record of intent that courts weigh when assessing voluntariness.

\subsection*{``I don’t consent to any searches.''}
This line asserts the Fourth Amendment’s protection against unreasonable searches and seizures. \textit{Schneckloth v.\ Bustamonte} (1973) ruled that police need not inform individuals of the right to refuse consent, and \textit{Florida v.\ Bostick} (1991) showed passive behavior can be construed as consent. An explicit refusal prevents claims of implied or voluntary consent.

\subsection*{``I want to leave. Am I free to go?''}
This line engages Fourth‑Amendment safeguards against unlawful detention. In \textit{Florida v.\ Royer} (1983) officers escalated a consensual conversation into an unlawful seizure without probable cause. \textit{Terry v.\ Ohio} (1968) permits brief investigative stops based on reasonable suspicion, but they must be limited in scope and duration. Saying ``I want to leave'' signals withdrawal of consent, while ``Am I free to go?'' forces clarification. Officers must then articulate reasonable suspicion to justify continued detention.

\bigskip
Each line is narrowly tailored for a discrete legal purpose—either a protection needing affirmative assertion or a clear evidentiary marker. Add‑ons like ``I reserve all rights'' have no clear legal effect and risk muddying the record; the Invocation includes only lines with functional weight in common encounters.

The script applies across all custody statuses and jurisdictions. Civilians need not diagnose encounter type in real time; the four lines work whether an interaction is consensual, detained, or arrested, drawing on cases such as \textit{Salinas} (pre‑custody), \textit{Berghuis} (post‑custody), \textit{Bostick} (consent searches), \textit{Schneckloth} (custodial consent), \textit{Royer} (consensual stops), and \textit{Terry} (brief detentions). States may expand protections but cannot diminish them, so the Invocation functions uniformly in all fifty states and territories.

Beyond immediate clarity, standardization lets courts and advocates assess effectiveness. If rulings reveal ambiguous phrasing, the script can be refined centrally and redistributed—a feedback loop impossible with improvised language. Officers also gain recognizable cues, simplifying proceedings by reducing disputes over intent.

The script respects legal limits. \textit{Hiibel v.\ Sixth Judicial District Court} (2004) upheld stop‑and‑identify statutes in some states, and traffic stops may require compliance with commands such as providing a license. Invocation documents intent; it does not override lawful orders.

The next section explains why memorized scripts survive stress and how behavioral science supports instinctive recall of legal language.
